\documentclass[11pt]{article}

\usepackage[margin=1in]{geometry}
\usepackage{amsmath,amssymb,amsthm,mathtools}
\usepackage{microtype}
\usepackage{enumitem}
\usepackage[colorlinks=true,linkcolor=blue,citecolor=blue,urlcolor=blue]{hyperref}

\newtheorem{theorem}{Theorem}[section]
\newtheorem{proposition}[theorem]{Proposition}
\newtheorem{lemma}[theorem]{Lemma}
\newtheorem{corollary}[theorem]{Corollary}
\newtheorem{remark}[theorem]{Remark}
\newtheorem{definition}[theorem]{Definition}

\newcommand{\R}{\mathbb{R}}
\newcommand{\one}{\mathbf{1}}
\newcommand{\ip}[2]{\left\langle #1,#2\right\rangle}
\newcommand{\Tr}{\operatorname{Tr}}
\newcommand{\Alt}{\operatorname{Alt}}
\newcommand{\op}{\mathrm{op}}
\newcommand{\HS}{\mathrm{HS}}

\title{Semi-inducibility of Alternating Cycles of Length
\texorpdfstring{$4k+2$}{4k+2}}
\author{Taeyoung Kim}
\date{\today}

\begin{document}
\maketitle

\begin{abstract}
Let $W\colon[0,1]^2\to[0,1]$ be a graphon, and let
$\Alt_{2m}(W)$ denote the density of the red--blue alternating cycle of
length $2m$, where red edges contribute $W$ and blue edges contribute
$1-W$.  We prove that, for every odd integer $m\ge 3$,
\[
  4^m\Alt_{2m}(W)+t(C_{2m},2W-1)\le 1.
\]
Since the signed even-cycle density $t(C_{2m},2W-1)$ is nonnegative,
this implies
\[
  \Alt_{2m}(W)\le 4^{-m}.
\]
Equivalently, for every $k\ge 1$, the alternating cycle of length
$4k+2$ has semi-induced density at most $2^{-(4k+2)}$.  Equality holds
only for the constant graphon $W\equiv 1/2$.  The proof reduces the
problem to an abstract rank-one matrix inequality.  A rank-two
Schur-complement factorisation produces a scalar power series whose
coefficients form a nonincreasing nonnegative sequence; since $m$ is
odd, comparing coefficients in the log-determinant expansion gives the
required sign.
\end{abstract}

\section{Definitions and main results}

A \emph{graphon} is a symmetric measurable function
$W\colon[0,1]^2\to[0,1]$.  We write
\[
  U:=1-W.
\]
For a bounded symmetric kernel $K\colon[0,1]^2\to\R$ and an integer
$r\ge 2$, its cycle density is
\[
  t(C_r,K)
  :=
  \int_{[0,1]^r}
  \prod_{j=1}^{r}K(x_j,x_{j+1})
  \,dx_1\cdots dx_r,
  \qquad x_{r+1}:=x_1.
\]
The same definition applies to the signed kernel $2W-1$.  For
background on graphons, homomorphism densities, and the cut norm, see
Lov\'asz~\cite{Lovasz}.

\begin{definition}[Alternating-cycle density]
For $m\ge 2$, define
\begin{align*}
  \Alt_{2m}(W)
  &:={}
  \int_{[0,1]^{2m}}
  \prod_{j=1}^{m}
  W(x_{2j-1},x_{2j})
  U(x_{2j},x_{2j+1})
  \,dx_1\cdots dx_{2m},
  \\
  &\hspace{8cm} x_{2m+1}:=x_1.
\end{align*}
Thus the colours alternate red, blue, red, blue around the cycle.
\end{definition}

\begin{theorem}[Strengthened alternating-cycle inequality]
\label{thm:main}
Let $m\ge 3$ be odd.  Then every graphon $W$ satisfies
\begin{equation}
  4^m\Alt_{2m}(W)+t(C_{2m},2W-1)\le 1.
  \label{eq:main-strengthened}
\end{equation}
Consequently,
\begin{equation}
  \Alt_{2m}(W)\le 4^{-m}.
  \label{eq:main-unweighted}
\end{equation}
Equality in \eqref{eq:main-unweighted} holds if and only if
$W=1/2$ almost everywhere.
\end{theorem}

Writing $m=2k+1$ gives the usual formulation.

\begin{corollary}[Alternating cycles of length $2\pmod 4$]
\label{cor:4k2}
For every integer $k\ge 1$ and every graphon $W$,
\[
  \Alt_{4k+2}(W)\le 2^{-(4k+2)}.
\]
Equality holds if and only if $W=1/2$ almost everywhere.
\end{corollary}

Alternating cycles belong to the semi-inducibility framework of Basit,
Granet, Horsley, K\"undgen, and Staden~\cite{BGHKS}, who proved sharp
or nearly sharp results for alternating cycles of length divisible by
four; the lengths congruent to $2$ modulo $4$ were left open.  The
first such case, length six, was settled by Chen and
Noel~\cite{ChenNoel} using flag algebras.
Corollary~\ref{cor:4k2} recovers that case ($k=1$) and answers the
question for the entire congruence class.

The strengthened inequality also gives a quantitative bound on
quasirandomness.  We use the cut norm
\[
  \|K\|_{\square}
  :=
  \sup_{S,T\subseteq[0,1]\ \mathrm{measurable}}
  \left|\int_{S\times T}K(x,y)\,dx\,dy\right|.
\]

\begin{corollary}[Stability]
\label{cor:stability}
Let $m\ge 3$ be odd and put
\[
  \delta:=1-4^m\Alt_{2m}(W).
\]
Then
\begin{equation}
  \left\|W-\frac12\right\|_{\square}
  \le
  \frac12\,\delta^{1/(2m)}.
  \label{eq:stability}
\end{equation}
\end{corollary}

We first reduce the graphon theorem to a finite-dimensional matrix
inequality, then prove that inequality by a Schur-complement expansion.

\section{Proof of the main results}

We include the approximation argument in full so that every determinant
appearing later is an ordinary finite-dimensional determinant.

Throughout this section we use the following dyadic approximation
scheme.  For $n\ge 1$, divide $[0,1]$ into the dyadic intervals
\[
  I_{n,a}:=
  \left[\frac{a}{2^n},\frac{a+1}{2^n}\right),
  \qquad 0\le a<2^n,
\]
with the convention that the last interval is closed, and let
$\mathcal A_n$ be the finite sigma-algebra generated by these
intervals.  For a bounded measurable kernel $K$ on $[0,1]^2$, set
\[
  K_n:=\mathbb E[K\mid \mathcal A_n\otimes\mathcal A_n].
\]
Then $K_n$ is a step kernel, constant on the rectangles
$I_{n,a}\times I_{n,b}$.  Conditional expectation preserves pointwise
bounds, so $\|K_n\|_\infty\le\|K\|_\infty$, and $0\le K\le1$ implies
$0\le K_n\le1$.  If $K$ is symmetric, then so is $K_n$, because
$\mathcal A_n\otimes\mathcal A_n$ is invariant under interchanging the
two coordinates.  Finally, the spaces of dyadic rectangle step
functions are nested and their union is dense in $L^2([0,1]^2)$; since
$K_n$ is the orthogonal projection of $K$ onto the $n$th such space,
\begin{equation}
  \|K_n-K\|_2\longrightarrow0,
  \qquad
  \|K_n-K\|_1\le\|K_n-K\|_2\longrightarrow0,
  \label{eq:dyadic-convergence}
\end{equation}
where the last inequality uses that $[0,1]^2$ has measure one.

\begin{lemma}[Step-graphon reduction]
\label{lem:step-reduction}
Fix $m\ge 2$.  To prove either
\eqref{eq:main-strengthened} or \eqref{eq:main-unweighted} for all
$W$, it is enough to prove it for symmetric step graphons.
\end{lemma}

\begin{proof}
Apply the dyadic scheme to $W$ itself: the kernels
\[
  W_n:=\mathbb E[W\mid \mathcal A_n\otimes\mathcal A_n]
\]
are symmetric step graphons with $0\le W_n\le1$, and
\eqref{eq:dyadic-convergence} gives $\|W_n-W\|_1\to0$.

We first check the alternating density.  For each point
$(x_1,\ldots,x_{2m})$, the integrand of $\Alt_{2m}(W_n)$ is a product
of $2m$ factors, each equal either to $W_n(x_i,x_{i+1})$ or to
$1-W_n(x_i,x_{i+1})$.  The corresponding factor for $W$ differs by
$\pm(W_n-W)(x_i,x_{i+1})$.  The telescoping identity for two products,
together with the fact that every factor lies in $[0,1]$, gives
\[
  |\Alt_{2m}(W_n)-\Alt_{2m}(W)|
  \le 2m\|W_n-W\|_1
  \longrightarrow0.
\]

Next put
\[
  K_n:=2W_n-1,
  \qquad
  K:=2W-1.
\]
All factors in the cycle integrands have absolute value at most $1$,
so the same telescoping argument gives
\[
  |t(C_{2m},K_n)-t(C_{2m},K)|
  \le 2m\|K_n-K\|_1
  =4m\|W_n-W\|_1
  \longrightarrow0.
\]
Thus every inequality proved for $W_n$ passes to $W$.
\end{proof}

We next describe the finite-dimensional operator model for step
kernels.  Let
\[
  [0,1]=B_1\sqcup\cdots\sqcup B_N
\]
be a finite measurable partition into sets of positive measure, and put
\[
  E:=\operatorname{span}\{\one_{B_1},\ldots,\one_{B_N}\}
  \subseteq L^2([0,1]).
\]
Every integral operator associated with a bounded kernel constant on
the rectangles $B_a\times B_b$ maps $L^2([0,1])$ into $E$ and vanishes
on $E^\perp$.  Consequently, all its nonzero action occurs on the
finite-dimensional real Hilbert space $E$.

For such a step kernel $K$, let $T_K$ be the integral operator
\[
  (T_Kf)(x):=\int_0^1K(x,y)f(y)\,dy.
\]
In the orthonormal basis
\[
  e_a:=\frac{\one_{B_a}}{\sqrt{|B_a|}},
\]
the matrix of $T_K|_E$ has entries
\[
  \ip{e_a}{T_Ke_b}=K_{ab}\sqrt{|B_a||B_b|}.
\]
It follows by direct multiplication that, for any sequence
$K_1,\ldots,K_r$ of step kernels constant on these rectangles,
\begin{equation}
  \Tr(T_{K_1}\cdots T_{K_r})
  =
  \int_{[0,1]^r}
  \prod_{j=1}^{r}K_j(x_j,x_{j+1})
  \,dx_1\cdots dx_r,
  \qquad x_{r+1}:=x_1.
  \label{eq:mixed-trace}
\end{equation}
In particular,
\begin{equation}
  \Tr(T_K^r)=t(C_r,K).
  \label{eq:cycle-trace}
\end{equation}

\begin{lemma}[Even signed cycles are nonnegative]
\label{lem:even-signed-cycle}
Let $K\colon[0,1]^2\to\R$ be a bounded symmetric measurable kernel,
and let $T_K$ be its integral operator on $L^2([0,1])$.  For every
integer $r\ge1$,
\begin{equation}
  t(C_{2r},K)
  =\Tr(T_K^{2r})
  =\|T_K^r\|_{\HS}^2
  \ge0.
  \label{eq:even-cycle-HS}
\end{equation}
\end{lemma}

\begin{proof}
First suppose that $K$ is a symmetric step kernel.  Then $T_K$ has
finite rank, the trace formula \eqref{eq:cycle-trace} gives
$t(C_{2r},K)=\Tr(T_K^{2r})$, and self-adjointness gives
\[
  \Tr(T_K^{2r})
  =\Tr\bigl((T_K^r)^*T_K^r\bigr)
  =\|T_K^r\|_{\HS}^2.
\]

For a general bounded symmetric kernel $K$, let
$K_n=\mathbb E[K\mid\mathcal A_n\otimes\mathcal A_n]$ be its dyadic
step approximations from the beginning of this section.  They are
symmetric, and \eqref{eq:dyadic-convergence} gives
\[
  \|K_n-K\|_2\to0,
  \qquad
  \|K_n-K\|_1\to0.
\]
Equivalently,
\[
  \|T_{K_n}-T_K\|_{\HS}\to0.
\]
The operators are uniformly bounded in operator norm, since
$\|T_{K_n}\|_{\op}\le\|T_{K_n}\|_{\HS}\le\|K\|_2$ and similarly for
$T_K$.  The telescoping identity
\[
  T_{K_n}^r-T_K^r
  =\sum_{j=0}^{r-1}
  T_{K_n}^{j}(T_{K_n}-T_K)T_K^{r-1-j}
\]
and the ideal inequality
$\|ABC\|_{\HS}\le\|A\|_{\op}\|B\|_{\HS}\|C\|_{\op}$ show that
$T_{K_n}^r\to T_K^r$ in Hilbert--Schmidt norm.  Hence
\[
  \|T_{K_n}^r\|_{\HS}^2\longrightarrow\|T_K^r\|_{\HS}^2.
\]
On the other hand, the same product-telescoping estimate as in
Lemma~\ref{lem:step-reduction}, now with
$M:=\|K\|_\infty$ and $\|K_n\|_\infty\le M$, gives
\[
  |t(C_{2r},K_n)-t(C_{2r},K)|
  \le 2rM^{2r-1}\|K_n-K\|_1
  \longrightarrow0.
\]
Passing to the limit in the step-kernel identity therefore gives
\[
  t(C_{2r},K)=\|T_K^r\|_{\HS}^2.
\]

Finally, $T_K$ is Hilbert--Schmidt and hence bounded, so the ideal
inequality gives
\[
  \|T_K^r\|_{\HS}
  \le\|T_K\|_{\op}^{r-1}\|T_K\|_{\HS}<\infty.
\]
Therefore $T_K^r$ is Hilbert--Schmidt.  Since $T_K$ is self-adjoint,
\[
  T_K^{2r}=(T_K^r)^*T_K^r
\]
is a positive trace-class operator, and
\[
  \Tr(T_K^{2r})=\|T_K^r\|_{\HS}^2.
\]
Combining the last two identities proves \eqref{eq:even-cycle-HS}.
\end{proof}

Now fix a symmetric step graphon $W$ and choose a partition as above
such that $W$ is constant on every rectangle $B_a\times B_b$.  The
constant function $e:=\one$ belongs to $E$ and has norm one.  Let
\[
  P:=e\otimes e,
\]
so that
\[
  Pf=\ip{e}{f}e.
\]
This is the rank-one orthogonal projection onto the constants, and it is
also the operator $T_1$ associated with the constant kernel $1$.
Finally set
\[
  X:=T_{2W-1}.
\]
Since $2W-1$ is symmetric, $X$ is self-adjoint.  Moreover,
\begin{equation}
  T_W=\frac{P+X}{2},
  \qquad
  T_U=\frac{P-X}{2}.
  \label{eq:color-coordinate}
\end{equation}

By \eqref{eq:mixed-trace},
\begin{equation}
  \Alt_{2m}(W)
  =\Tr\bigl((T_WT_U)^m\bigr)
  =4^{-m}\Tr\left(\bigl((P+X)(P-X)\bigr)^m\right).
  \label{eq:alt-trace}
\end{equation}
Also,
\begin{equation}
  \Tr(X^2)
  =\int_{[0,1]^2}(2W(x,y)-1)^2\,dx\,dy
  \le1.
  \label{eq:HS-budget}
\end{equation}
Thus Theorem~\ref{thm:main} follows from the following abstract matrix
theorem.

\begin{theorem}[Abstract alternating trace inequality]
\label{thm:matrix}
Let $H$ be a finite-dimensional real Hilbert space, let $e\in H$ be a
unit vector, and let
\[
  P=e\otimes e.
\]
Let $X=X^*$ satisfy
\[
  \tau:=\Tr(X^2)\le1.
\]
For every odd positive integer $m$,
\begin{equation}
  \Tr\left(\bigl((P+X)(P-X)\bigr)^m\right)
  +\Tr(X^{2m})
  \le1.
  \label{eq:matrix-main}
\end{equation}
\end{theorem}

The proof is given in the remainder of this section.  Put
\begin{equation}
  L:=(P+X)(P-X).
  \label{eq:def-L}
\end{equation}
Although $L$ need not be self-adjoint, the determinant--trace identity
used below is valid for every finite matrix.

All power-series identities below may be read either formally
or as Taylor expansions for $z$ sufficiently close to zero.  Since all
matrices are finite-dimensional, the two interpretations agree
coefficient by coefficient.

Define
\[
  D(z):=I+zX^2,
  \qquad
  N(z):=D(z)^{-1},
  \qquad
  u:=Xe.
\]
Also define the scalar functions
\begin{align}
  h(z)&:=\ip{e}{N(z)e},
  \label{eq:def-h}\\
  k(z)&:=\ip{e}{XN(z)e}=\ip{u}{N(z)e},
  \label{eq:def-k}\\
  \ell(z)&:=\ip{e}{X^2N(z)e}=\ip{u}{N(z)u}.
  \label{eq:def-ell}
\end{align}
Since
\[
  N(z)+zX^2N(z)=I,
\]
we have
\begin{equation}
  h(z)+z\ell(z)=1.
  \label{eq:h-ell}
\end{equation}

\begin{lemma}[Rank-two determinant factorisation]
\label{lem:det-factor}
With $L$ as in \eqref{eq:def-L},
\begin{equation}
  \det(I-zL)
  =
  \det(I+zX^2)
  \bigl(1-zF(z)\bigr),
  \label{eq:det-factor}
\end{equation}
where
\begin{equation}
  F(z):=h(z)^2+zk(z)^2.
  \label{eq:def-F}
\end{equation}
\end{lemma}

\begin{proof}
Expanding $L$ gives
\[
  L=P-PX+XP-X^2.
\]
Since $PX=e\otimes u$ and $XP=u\otimes e$,
\begin{align*}
  I-zL
  &=I+zX^2-zP+zPX-zXP\\
  &=D(z)+z\,e\otimes(u-e)-z\,u\otimes e.
\end{align*}
Introduce the maps $\mathcal U\colon\R^2\to H$ and
$\mathcal V\colon\R^2\to H$ whose columns are
\[
  \mathcal U=\begin{bmatrix}ze&-zu\end{bmatrix},
  \qquad
  \mathcal V=\begin{bmatrix}u-e&e\end{bmatrix}.
\]
Then
\[
  I-zL=D(z)+\mathcal U\mathcal V^*.
\]
The matrix determinant lemma gives
\[
  \det(I-zL)
  =\det D(z)\,
  \det\bigl(I_2+\mathcal V^*N(z)\mathcal U\bigr).
\]
Because $N(z)$ is a function of $X^2$, it commutes with $X$.  Using
\eqref{eq:def-h}--\eqref{eq:def-ell}, we obtain
\[
  I_2+\mathcal V^*N(z)\mathcal U
  =
  \begin{pmatrix}
    1+z(k-h) & z(k-\ell)\\
    zh       & 1-zk
  \end{pmatrix}.
\]
Its determinant is
\begin{align*}
 &\bigl(1+z(k-h)\bigr)(1-zk)-z^2h(k-\ell)\\
 &\qquad
 =1-zh-z^2k^2+z^2h\ell.
\end{align*}
Using \eqref{eq:h-ell}, namely $z\ell=1-h$, this becomes
\[
  1-zh^2-z^2k^2
  =1-z\bigl(h^2+zk^2\bigr).
\]
Substitution proves \eqref{eq:det-factor}.
\end{proof}

Diagonalise $X$ in an orthonormal eigenbasis:
\[
  Xv_i=\lambda_i v_i,
  \qquad
  \omega_i:=|\ip{e}{v_i}|^2.
\]
Eigenvalues are listed with multiplicity.  Then
\begin{equation}
  \omega_i\ge0,
  \qquad
  \sum_i\omega_i=1,
  \qquad
  \sum_i\lambda_i^2=\tau\le1.
  \label{eq:spectral-budget}
\end{equation}
The scalar resolvents are
\begin{equation}
  h(z)=\sum_i\frac{\omega_i}{1+z\lambda_i^2},
  \qquad
  k(z)=\sum_i\frac{\omega_i\lambda_i}{1+z\lambda_i^2}.
  \label{eq:h-k-spectral}
\end{equation}
Consequently,
\begin{equation}
  F(z)
  =
  \sum_{i,j}\omega_i\omega_j
  \frac{1+z\lambda_i\lambda_j}
       {(1+z\lambda_i^2)(1+z\lambda_j^2)}.
  \label{eq:F-double}
\end{equation}

For $n\ge0$, define the polynomial
\begin{equation}
  c_n(x,y)
  :=
  \frac{x^{2n+1}+y^{2n+1}}{x+y},
  \label{eq:def-cn-quotient}
\end{equation}
where the quotient is interpreted by continuity at $x+y=0$.  Equivalently,
\begin{equation}
  c_n(x,y)
  =
  \sum_{r=0}^{2n}(-1)^r x^{2n-r}y^r.
  \label{eq:def-cn-polynomial}
\end{equation}
In particular,
\[
  c_n(x,-x)=(2n+1)x^{2n}.
\]

\begin{lemma}[Positivity and recurrence for $c_n$]
\label{lem:cn}
For all real $x,y$ and all $n\ge0$,
\begin{equation}
  c_n(x,y)\ge0.
  \label{eq:cn-positive}
\end{equation}
Moreover, for $n\ge1$,
\begin{equation}
  (x^2+y^2)c_n(x,y)-c_{n+1}(x,y)
  =x^2y^2c_{n-1}(x,y).
  \label{eq:cn-recurrence}
\end{equation}
\end{lemma}

\begin{proof}
Let $f(t)=t^{2n+1}$.  If $x+y\ne0$, then
\[
  c_n(x,y)=\frac{f(x)-f(-y)}{x-(-y)}.
\]
The function $f$ is increasing on $\R$, so this difference quotient,
the slope of a chord of $f$, is nonnegative.  If $x=-y$, the continuous value is
$f'(x)=(2n+1)x^{2n}\ge0$.  This proves
\eqref{eq:cn-positive}.

For the recurrence, multiply the claimed identity by $x+y$.  The left
side becomes
\begin{align*}
 &(x^2+y^2)(x^{2n+1}+y^{2n+1})
 -(x^{2n+3}+y^{2n+3})\\
 &\qquad
 =x^2y^{2n+1}+y^2x^{2n+1}\\
 &\qquad
 =x^2y^2(x^{2n-1}+y^{2n-1}),
\end{align*}
which is $(x+y)x^2y^2c_{n-1}(x,y)$.  Since both sides are polynomials,
the identity holds also when $x+y=0$.
\end{proof}

The partial-fraction identity
\begin{equation}
  \frac{1+zxy}{(1+zx^2)(1+zy^2)}
  =
  \frac{x}{x+y}\frac1{1+zx^2}
  +
  \frac{y}{x+y}\frac1{1+zy^2}
  \label{eq:partial-fraction}
\end{equation}
for $x+y\ne0$, extended continuously to $x+y=0$, gives
\begin{equation}
  \frac{1+zxy}{(1+zx^2)(1+zy^2)}
  =
  \sum_{n\ge0}(-1)^n c_n(x,y)z^n.
  \label{eq:cn-generating}
\end{equation}
Combining \eqref{eq:F-double} and \eqref{eq:cn-generating}, write
\begin{equation}
  F(z)=\sum_{n\ge0}(-1)^n\beta_nz^n,
  \label{eq:F-beta}
\end{equation}
where
\begin{equation}
  \beta_n:=\sum_{i,j}\omega_i\omega_j
  c_n(\lambda_i,\lambda_j).
  \label{eq:def-beta}
\end{equation}

\begin{lemma}[Monotonicity of the coefficients $\beta_n$]
\label{lem:beta-monotone}
The sequence $(\beta_n)_{n\ge0}$ satisfies
\begin{equation}
  1=\beta_0\ge\beta_1\ge\beta_2\ge\cdots\ge0.
  \label{eq:beta-monotone}
\end{equation}
More precisely,
\begin{equation}
  \beta_{n+1}\le\tau\beta_n
  \qquad(n\ge1).
  \label{eq:beta-tau-contract}
\end{equation}
\end{lemma}

\begin{proof}
Positivity follows from Lemma~\ref{lem:cn}.  Since $c_0\equiv1$ and
$\sum_i\omega_i=1$, we have $\beta_0=1$.

We first prove $\beta_1\le\tau$.  Put
\[
  \mu_r:=\sum_i\omega_i\lambda_i^r=\ip{e}{X^re}.
\]
Since
\[
  c_1(x,y)=x^2-xy+y^2,
\]
we obtain
\begin{equation}
  \beta_1=2\mu_2-\mu_1^2.
  \label{eq:beta1-moments}
\end{equation}
Decompose $H=\R e\oplus e^\perp$ and write
\[
  X=
  \begin{pmatrix}
    a&w^*\\
    w&A
  \end{pmatrix}.
\]
Then
\[
  \mu_1=a,
  \qquad
  \mu_2=a^2+\|w\|^2,
\]
and hence
\begin{equation}
  \beta_1=a^2+2\|w\|^2.
  \label{eq:beta1-block}
\end{equation}
On the other hand,
\begin{equation}
  \tau=\Tr(X^2)=a^2+2\|w\|^2+\Tr(A^2).
  \label{eq:tau-block}
\end{equation}
Thus $\beta_1\le\tau\le1$.

Now let $n\ge1$.  For a diagonal pair $i=j$, one has
\[
  c_n(\lambda_i,\lambda_i)=\lambda_i^{2n},
\]
so
\[
  c_{n+1}(\lambda_i,\lambda_i)
  =\lambda_i^2c_n(\lambda_i,\lambda_i)
  \le\tau c_n(\lambda_i,\lambda_i).
\]
For $i\ne j$, Lemma~\ref{lem:cn} gives
\begin{align*}
  c_{n+1}(\lambda_i,\lambda_j)
  &\le
  (\lambda_i^2+\lambda_j^2)c_n(\lambda_i,\lambda_j)\\
  &\le
  \tau c_n(\lambda_i,\lambda_j),
\end{align*}
where the last inequality follows because the eigenvalues are listed
with multiplicity and
\[
  \lambda_i^2+\lambda_j^2\le\sum_s\lambda_s^2=\tau.
\]
Averaging with the nonnegative weights $\omega_i\omega_j$ proves
\eqref{eq:beta-tau-contract}.  Since $\tau\le1$, the sequence is
nonincreasing.
\end{proof}

The following purely one-dimensional lemma isolates the role of the
parity of $m$.

\begin{lemma}[Odd logarithmic coefficient]
\label{lem:odd-log}
Let
\[
  1=\beta_0\ge\beta_1\ge\beta_2\ge\cdots
\]
and put
\[
  F(z):=\sum_{n\ge0}(-1)^n\beta_nz^n.
\]
Then, for every odd positive integer $m$,
\begin{equation}
  m[z^m]\bigl(-\log(1-zF(z))\bigr)\le1.
  \label{eq:odd-log-bound}
\end{equation}
More precisely, if
\[
  d_n:=\beta_n-\beta_{n+1}\ge0,
  \qquad
  G(z):=\sum_{n\ge0}d_nz^n,
\]
then
\begin{equation}
  1-m[z^m]\bigl(-\log(1-zF(z))\bigr)
  =
  m\sum_{r=1}^{(m-1)/2}\frac1r
  [z^{m-2r}]G(z)^r.
  \label{eq:odd-log-defect}
\end{equation}
In particular, the right side is nonnegative.
\end{lemma}

\begin{proof}
The coefficient of $z^n$ in $(1+z)F(z)$ is the successive difference
of the $\beta_j$ with alternating sign.  Explicitly,
\begin{equation}
  (1+z)F(z)=1+zG(-z).
  \label{eq:F-G}
\end{equation}
Therefore
\begin{align}
  (1+z)(1-zF(z))
  &=1+z-z(1+z)F(z)\notag\\
  &=1-z^2G(-z).
  \label{eq:factor-log}
\end{align}
Taking formal logarithms gives
\begin{equation}
  -\log(1-zF(z))
  =
  \log(1+z)-\log(1-z^2G(-z)).
  \label{eq:log-decomposition}
\end{equation}
For odd $m$,
\[
  [z^m]\log(1+z)=\frac1m.
\]
Furthermore,
\begin{equation}
  -\log(1-z^2G(-z))
  =
  \sum_{r\ge1}\frac{z^{2r}G(-z)^r}{r}.
  \label{eq:log-G-expansion}
\end{equation}
If $m$ is odd, then $m-2r$ is odd.  Since $G$ has nonnegative
coefficients,
\[
  [z^{m-2r}]G(-z)^r
  =-[z^{m-2r}]G(z)^r\le0.
\]
Substituting this into \eqref{eq:log-decomposition} gives the exact
identity \eqref{eq:odd-log-defect}, and hence
\eqref{eq:odd-log-bound}.
\end{proof}

We use the standard finite-dimensional identity
\begin{equation}
  -\log\det(I-zM)
  =
  \sum_{r\ge1}\frac{\Tr(M^r)}{r}z^r
  \label{eq:logdet-trace}
\end{equation}
for an arbitrary finite matrix $M$.  For completeness, differentiating
the left side and using Jacobi's formula gives
\[
  \frac{d}{dz}\bigl[-\log\det(I-zM)\bigr]
  =\Tr\bigl(M(I-zM)^{-1}\bigr)
  =\sum_{r\ge0}\Tr(M^{r+1})z^r.
\]
Both sides of \eqref{eq:logdet-trace} vanish at $z=0$, proving the
identity coefficientwise.  No diagonalizability of $M$ is required.

\begin{proof}[Proof of Theorem~\ref{thm:matrix}]
Take minus the logarithm of the determinant factorisation
\eqref{eq:det-factor}:
\begin{equation}
  -\log\det(I-zL)
  =
  -\log\det(I+zX^2)
  -\log(1-zF(z)).
  \label{eq:logdet-factor}
\end{equation}
By \eqref{eq:logdet-trace},
\[
  -\log\det(I-zL)
  =\sum_{r\ge1}\frac{\Tr(L^r)}{r}z^r.
\]
Since $X$ is self-adjoint,
\begin{equation}
  -\log\det(I+zX^2)
  =
  \sum_{r\ge1}\frac{(-1)^r\Tr(X^{2r})}{r}z^r.
  \label{eq:X-even-log}
\end{equation}
Let $m$ be odd.  Comparing the coefficient of $z^m$ in
\eqref{eq:logdet-factor} gives
\begin{equation}
  \Tr(L^m)
  =
  -\Tr(X^{2m})
  +m[z^m]\bigl(-\log(1-zF(z))\bigr).
  \label{eq:trace-coefficient}
\end{equation}
By Lemma~\ref{lem:beta-monotone}, the coefficients of $F$ satisfy the
hypotheses of Lemma~\ref{lem:odd-log}.  Hence
\[
  m[z^m]\bigl(-\log(1-zF(z))\bigr)\le1.
\]
Substitution into \eqref{eq:trace-coefficient} proves
\[
  \Tr(L^m)+\Tr(X^{2m})\le1.
\]
This is \eqref{eq:matrix-main}.
\end{proof}

The proof actually supplies an exact formula for the gap.

\begin{proposition}[Exact formula for the gap]
\label{prop:matrix-defect}
With the notation above, let
\[
  d_n:=\beta_n-\beta_{n+1},
  \qquad
  G(z):=\sum_{n\ge0}d_nz^n.
\]
For every odd positive integer $m$,
\begin{equation}
  1-\Tr(L^m)-\Tr(X^{2m})
  =
  m\sum_{r=1}^{(m-1)/2}\frac1r
  [z^{m-2r}]G(z)^r.
  \label{eq:matrix-defect}
\end{equation}
Every term on the right is nonnegative.
\end{proposition}

\begin{proof}
Combine \eqref{eq:trace-coefficient} with the exact identity
\eqref{eq:odd-log-defect}.
\end{proof}

\begin{proof}[Proof of Theorem~\ref{thm:main}]
By Lemma~\ref{lem:step-reduction}, first assume that $W$ is a symmetric
step graphon.  Work on the finite-dimensional space $E$ from the
finite-rank reduction, and define $P$ and $X$ as in
\eqref{eq:color-coordinate}.
The bound \eqref{eq:HS-budget} gives
\[
  \Tr(X^2)\le1.
\]
Therefore Theorem~\ref{thm:matrix}, together with
\eqref{eq:alt-trace} and \eqref{eq:cycle-trace}, yields
\begin{align*}
  4^m\Alt_{2m}(W)+t(C_{2m},2W-1)
  &=
  \Tr\left(\bigl((P+X)(P-X)\bigr)^m\right)
  +\Tr(X^{2m})\\
  &\le1.
\end{align*}
This proves \eqref{eq:main-strengthened} for step graphons.  The
approximation in Lemma~\ref{lem:step-reduction} passes it to every
graphon.

Lemma~\ref{lem:even-signed-cycle}, applied to the signed kernel
$2W-1$, gives
\[
  t(C_{2m},2W-1)=\Tr(X^{2m})=\|X^m\|_{\HS}^2\ge0.
\]
Dropping this nonnegative term from
\eqref{eq:main-strengthened} proves \eqref{eq:main-unweighted}.

Suppose equality holds in \eqref{eq:main-unweighted}.  Then
\eqref{eq:main-strengthened} and Lemma~\ref{lem:even-signed-cycle} give
\[
  0\le\Tr(X^{2m})
  =\|X^m\|_{\HS}^2
  \le1-4^m\Alt_{2m}(W)
  =0.
\]
Hence $X^m=0$.  Because $X$ is self-adjoint,
\[
  \|X\|_{\op}^m=\|X^m\|_{\op}=0,
\]
so $X=0$.  The Hilbert--Schmidt kernel of $X$ is $2W-1$, and therefore
$2W-1=0$ almost everywhere.  Hence $W=1/2$ almost everywhere.
Conversely, $W\equiv1/2$ gives
\[
  \Alt_{2m}(W)=2^{-2m}=4^{-m},
\]
so equality does occur.
\end{proof}

\begin{proof}[Proof of Corollary~\ref{cor:stability}]
Put $X=T_{2W-1}$.  The strengthened inequality gives
\[
  \Tr(X^{2m})\le\delta.
\]
By \eqref{eq:main-unweighted}, $\delta\ge0$.  If $(\lambda_j)$ are the
eigenvalues of the compact self-adjoint operator $X$, counted with
multiplicity, then
\[
  \|X\|_{\op}^{2m}
  =\sup_j|\lambda_j|^{2m}
  \le\sum_j|\lambda_j|^{2m}
  =\Tr(X^{2m})
  \le\delta.
\]
For measurable $S,T\subseteq[0,1]$,
\[
  \left|\int_{S\times T}(2W-1)\right|
  =|\ip{\one_S}{X\one_T}|
  \le\|X\|_{\op}\|\one_S\|_2\|\one_T\|_2
  \le\|X\|_{\op}.
\]
Hence
\[
  \|2W-1\|_{\square}\le\|X\|_{\op}\le\delta^{1/(2m)}.
\]
Dividing by $2$ proves \eqref{eq:stability}.
\end{proof}

\section{Remarks}

Although the theorem is an unrestricted semi-inducibility result, the
signed-cycle term also gives a simple density-sensitive bound.

\begin{corollary}[A fixed-edge-density upper bound]
\label{cor:fixed-density}
Let
\[
  p:=\int_{[0,1]^2}W(x,y)\,dx\,dy.
\]
For every odd $m\ge3$,
\begin{equation}
  \Alt_{2m}(W)
  \le
  4^{-m}\bigl(1-(2p-1)^{2m}\bigr).
  \label{eq:fixed-density}
\end{equation}
\end{corollary}

\begin{proof}
For $X=T_{2W-1}$ and $e=\one$,
\[
  \ip{e}{Xe}=2p-1.
\]
Let $\nu$ be the spectral probability measure of $X$ at $e$.  Since
$t\mapsto t^{2m}$ is convex,
\[
  \ip{e}{X^{2m}e}
  =\int t^{2m}\,d\nu(t)
  \ge
  \left(\int t\,d\nu(t)\right)^{2m}
  =(2p-1)^{2m}.
\]
Also $X^{2m}$ is positive semidefinite, so
\[
  \Tr(X^{2m})\ge\ip{e}{X^{2m}e}.
\]
Substitute this into \eqref{eq:main-strengthened}.
\end{proof}

The inequalities \eqref{eq:main-strengthened} and
\eqref{eq:main-unweighted}, together with the fact that $W\equiv1/2$
attains the equality, have been formalised in the Lean~4 proof
assistant~\cite{Lean4}, on top of the mathlib
library~\cite{mathlib}; the development accompanies this note.
Corollaries~\ref{cor:stability} and~\ref{cor:fixed-density} and the
uniqueness of the equality case are not formalised.  The formal proof
differs from the one given here in three ways.

First, instead of the step-graphon approximation
(Lemma~\ref{lem:step-reduction}), we write the graphon quantity
$4^m\Alt_{2m}(W)+t(C_{2m},2W-1)$ and its matrix counterpart as the
same polynomial in the moments $\mu_g=\ip{\one}{X^g\one}$, $g<2m$,
and compress $X$ to the Krylov subspace
$\operatorname{span}\{\one,X\one,\ldots,X^{2m}\one\}$; this produces
a finite symmetric matrix with the same moments and $\Tr(A^2)\le1$,
to which Theorem~\ref{thm:matrix} applies directly.

Second, the formal proof never forms a logarithm or a determinant.
Instead of handling $m[z^m]\bigl(-\log A(z)\bigr)$, it uses the
operator $\Lambda(A):=-zA'(z)/A(z)$: since $z\,d/dz$ multiplies the
coefficient of $z^m$ by $m$, one has
$[z^m]\Lambda(A)=m[z^m]\bigl(-\log A(z)\bigr)$, while $\Lambda$
itself involves no logarithm.  The rule
$\Lambda(fg)=\Lambda(f)+\Lambda(g)$, immediate from
$(fg)'=f'g+fg'$, replaces the only property of the logarithm that
the argument uses.  Similarly, the logarithm in
\eqref{eq:trace-coefficient} is read as $[z^m]\Lambda(1-zF(z))$, and
the identity is proved not through \eqref{eq:logdet-trace} but
directly from the rank-two structure in the proof of
Lemma~\ref{lem:det-factor}, as a computation with the $2\times2$
matrix $I_2+\mathcal V^*N(z)\mathcal U$; the only determinants that
appear are $2\times2$.

Third, in Lemma~\ref{lem:beta-monotone}, we obtain $\beta_1\le\tau$
from a direct sum-of-squares identity, without decomposing $H$ along
$e$.

\begin{thebibliography}{9}

\bibitem{BGHKS}
A. Basit, B. Granet, D. Horsley, A. K\"undgen, and K. Staden,
\emph{The semi-inducibility problem},
arXiv:2501.09842.

\bibitem{ChenNoel}
H. Chen and J. A. Noel,
\emph{On alternating 6-cycles in edge-coloured graphs},
arXiv:2505.09809.

\bibitem{Lovasz}
L. Lov\'asz,
\emph{Large Networks and Graph Limits},
American Mathematical Society Colloquium Publications, vol.~60,
American Mathematical Society, Providence, RI, 2012.

\bibitem{mathlib}
The mathlib Community,
\emph{The Lean mathematical library},
Proceedings of the 9th ACM SIGPLAN International Conference on
Certified Programs and Proofs (CPP 2020), ACM, 2020, pp.~367--381.

\bibitem{Lean4}
L. de Moura and S. Ullrich,
\emph{The Lean 4 theorem prover and programming language},
Automated Deduction -- CADE 28, Lecture Notes in Computer Science,
vol.~12699, Springer, 2021, pp.~625--635.

\end{thebibliography}

\end{document}
